\section{Methodology}
\label{sec:method}

\begin{figure*}[t]
  \centering
  % Method overview (Fig. 1). Included from sections/methodology.tex.
\begin{tikzpicture}[
    font=\footnotesize,
    >={Stealth[length=4pt]},
    box/.style={draw, rounded corners=2pt, align=center, minimum height=8mm,
                inner sep=3pt, fill=#1},
    box/.default=white,
    lbl/.style={font=\scriptsize, align=center},
  ]
  % Row 1: input, ISP, condition branch
  \node[box=gray!10, text width=13mm] (raw) at (0,0) {Short-exp.\\raw};
  \node[box=gray!10, text width=14mm, right=5mm of raw] (isp) {Minimal\\ISP};
  \node[box=gray!10, text width=13mm, right=5mm of isp] (lrgb) {Noisy\\LRGB};
  \node[box=blue!8, text width=15mm, right=6mm of lrgb] (enc) {VAE enc.\ $\Enc$\\(frozen)};
  \node[box=orange!15, text width=19mm, right=9mm of enc] (ctx) {Context processor\\(U-Net enc.\ copy)};

  % Row 2: illumination, denoiser, decoder, output
  \node[box=yellow!25, text width=17mm, below=11mm of lrgb] (ill) {Illum.\ map $\illum$\\$\phi(\illum)=[\illum,\log \illum]$};
  \node[box=orange!15, text width=26mm, below=11mm of ctx] (unet) {\textbf{Compact U-Net}\\region cross-attn\\+ FiLM, key bias, gate};
  \node[box=blue!8, text width=17mm, right=9mm of unet] (dec) {Residual dec.\ $\Dec$\\(frozen)};
  \node[box=gray!10, text width=11mm, right=5mm of dec] (out) {sRGB\\output};
  \node[lbl, below=6mm of unet] (zt) {$z_T \sim \mathcal{N}(0, I)$, 50 DDIM steps, CFG};

  \draw[->] (raw) -- (isp);
  \draw[->] (isp) -- (lrgb);
  \draw[->] (lrgb) -- (enc);
  \draw[->] (enc) -- node[lbl, above] {$z_y$} (ctx);
  \draw[->] (ctx) -- node[lbl, right] {context tokens} (unet);
  \draw[->] (lrgb) -- node[lbl, left] {luminance,\\avg-pool} (ill);
  \draw[->, thick, yellow!50!black] (ill) -- (unet);
  \draw[->] (zt) -- (unet);
  \draw[->] (unet) -- node[lbl, above] {$\hat z_0$} (dec);
  \draw[->] (dec) -- (out);
  % Encoder-to-decoder skip features of the residual VAE
  \draw[->, dashed] (enc.north) -- ++(0,4mm) -| node[lbl, pos=0.25, above] {encoder skip features $E_b$} (dec.north);
\end{tikzpicture}

  \caption{Overview of our model. The DarkDiff pipeline~\cite{zheng2025darkdiff} is kept, but the
  \sdtwo{} U-Net is replaced by a distilled compact U-Net, and an illumination map of the input
  (yellow) conditions every region cross-attention layer. Blue blocks are frozen; orange
  blocks are trained.}
  \label{fig:pipeline}
\end{figure*}

Fig.~\ref{fig:pipeline} shows the full model. We first describe our re-implementation of
DarkDiff, then the two extensions.

\subsection{DarkDiff Pipeline}
\label{sec:darkdiff}

\textbf{Input.}
A minimal ISP packs the Bayer mosaic into four half-resolution planes, subtracts the black level,
normalises by the white level, multiplies the short exposure by the exposure ratio, and applies
white-balance gains. Averaging the two green planes gives a noisy linear RGB (LRGB) image $y$.
The long-exposure reference goes through the same steps, then a colour-correction matrix from the
camera metadata and the sRGB gamma curve, giving the target $x_{\mathrm{sRGB}}$.

\textbf{Residual VAE.}
The Stable Diffusion VAE downsamples by $8\times$, which loses fine detail. DarkDiff adds a
convolution from each encoder stage $E_b$ to the matching decoder stage, so the decoder sees
high-resolution features of the input. We initialise these skip convolutions at zero, so the VAE
starts as the pretrained one. The VAE is frozen while the diffusion model trains.

\textbf{Context processor and region cross-attention.}
The condition latent is $z_y = \Enc(y)$. A copy of the U-Net encoder (the context processor)
turns $z_y$ into features $c$ at every resolution. It runs at timestep 0, so it is computed once
per image rather than once per denoising step. In the down and mid blocks of the U-Net, the text
cross-attention is replaced by \emph{region} cross-attention. The token grid is split into
$8\times 8$ windows, and the queries $Q_w$ of window $w$ attend only to the context tokens of the
same window:
\begin{equation}
  \mathrm{RA}(Q_w, K_w, V_w) = \mathrm{softmax}\!\left(\frac{Q_w K_w^{\top}}{\sqrt{d}} + B\right) V_w ,
  \label{eq:region}
\end{equation}
where $K_w$ and $V_w$ are projections of $c$, and $B=0$ in plain DarkDiff. A pixel of a
processed image depends mostly on its neighbourhood, so this restriction keeps the output faithful
to the input's structure and limits hallucination. The projections are copied from the block's
self-attention, and the output projection starts at zero.

\textbf{Training loss.}
The U-Net $f_\theta$ predicts the noise $\epsilon$ or the velocity $v$~\cite{salimans2022progressive}
from the noisy latent $z_t$. From its prediction we compute a one-step estimate $\hat z_0$ of the
clean latent, decode it with the frozen residual decoder and compare it with the sRGB target:
\begin{equation}
  \mathcal{L} = \Lldm + \lambda\,\Limg, \qquad
  \Limg = \big\| \Dec(\hat z_0, E_b) - x_{\mathrm{sRGB}} \big\|_2^2 ,
  \label{eq:loss}
\end{equation}
with $\lambda = 1$. The image loss curbs colour shifts that the latent loss alone does not see.

\textbf{Guidance.}
During training, 5\% of the condition latents are replaced by Gaussian noise. At inference, the
prediction is $\hat f = f_u + w\,(f_c - f_u)$, where $f_c$ uses the real condition and $f_u$ a
noise condition. We use 50 DDIM steps and $w = 2.0$ for SID.

\subsection{Compact Backbones}
\label{sec:backbones}

The context processor is a copy of the U-Net encoder, so a smaller U-Net shrinks both the
denoiser and the context processor. The new U-Net must be \emph{pretrained} in the Stable
Diffusion latent space, because DarkDiff's quality comes from the pretrained prior. In early
tests, two compact U-Nets we trained from random initialisation (a half-width SD U-Net with
150M parameters and a MobileNetV3-style~\cite{howard2019mobilenetv3} variant with 100M) gave
unusable images after 10 epochs. We therefore use the three distilled U-Nets in
Table~\ref{tab:backbones}, which are the smallest public SD-compatible U-Nets we found.

\begin{table}[t]
  \centering
  \caption{Diffusion backbones. Sizes count U-Net parameters only.}
  \label{tab:backbones}
  \setlength{\tabcolsep}{4pt}
  \begin{tabular}{@{}llrc@{}}
    \toprule
    Backbone & Distilled from & U-Net & Pred. \\
    \midrule
    \sdtwo{} (baseline)            & ---                                & 866M & $v$ \\
    \bksdm{}~\cite{kim2024bksdm}   & SD 2.1-base                        & 329M & $\epsilon$ \\
    \tinysd{}~\cite{segmind2023tinysd} & Realistic Vision 4.0 (SD 1.5)  & 327M & $\epsilon$ \\
    \sdxs{}~\cite{song2024sdxs}    & SD 2.1-base, via SD-Turbo          & 330M & $\epsilon$ \\
    \bottomrule
  \end{tabular}
\end{table}

All three remove one residual and attention pair per stage, the mid block, and the deepest
1280-channel resolution level. The remaining layers keep SD's widths, so they are initialised
from the teacher before distillation. Most parameters live in the deep, low-resolution blocks,
while most FLOPs come from the high-resolution blocks that are kept. Compute therefore drops less
than size: 62\% fewer parameters, but only 39--47\% fewer FLOPs (Table~\ref{tab:efficiency}).
\sdxs{} replaces the self-attention of its highest-resolution blocks by text cross-attention.
There we copy only the projection weights whose shapes match.

To make the larger models trainable on one 48\,GB GPU, we use gradient checkpointing in the
VAE decoder. The backward pass of $\Limg$ through the full-resolution decoder dominated memory,
and checkpointing cut its activation memory by about 58\%.

\subsection{Illumination-Aware Conditioning}
\label{sec:illum}

Dark regions of the input have the lowest SNR, so the model should rely less on them and more
on its prior. We give every region cross-attention layer an illumination map that tells it how
bright each region of the input is.

\textbf{Illumination map.}
We compute the luminance $\ell = 0.2126R + 0.7152G + 0.0722B$ of the gain-amplified noisy LRGB,
average-pool it to the latent grid, clamp it to $[0,1]$ and smooth it with a $3\times3$ filter:
\begin{equation}
  \illum = \mathrm{smooth}\big(\mathrm{clamp}(\mathrm{avgpool}(\ell), 0, 1)\big).
\end{equation}
Pooling over $8\times8$ pixels before clamping reduces the noise by about $8\times$, so clipping
zero-mean noise does not bias the estimate upwards. The map is resized to each layer's token grid.
We use the features
$\phi(\illum) = \big[\illum,\ \log(\illum+\varepsilon)/\log(1/\varepsilon)\big]$ with
$\varepsilon = 10^{-3}$. The log channel turns exposure ratios into differences.

\textbf{Conditioning.}
The map enters each region cross-attention layer in three places.
\begin{enumerate}
  \item \emph{FiLM on the condition}, before the key and value projections:
        $c' = c \odot \big(1 + \gamma(\phiL)\big) + \beta(\phiL)$, where $\gamma$ and $\beta$
        are two-layer convolutions.
  \item \emph{Key bias}: in~\eqref{eq:region}, $B_{h,k} = b_h \log(\illum_k + \varepsilon)$ for
        head $h$ and key $k$, with a learned scalar $b_h$. A head can learn to down-weight keys
        from very dark, noisy tokens.
  \item \emph{Output gate}: the attention output is multiplied by
        $g = 1 + \tanh\big(G(\phiL)\big) \in (0, 2)$, so each region decides how much of the
        condition to inject.
\end{enumerate}
The last layers of $\gamma$, $\beta$ and $G$, and every $b_h$, are initialised to zero. At the
first step the model is therefore exactly DarkDiff. When the condition is replaced by noise
(in training or in the unconditional CFG branch), we set $\illum = 1$, which is neutral because
$\log \illum = 0$. The extension adds about 2M parameters and 2--3\% FLOPs.
